\section{Related Work}\label{sec:related}

\paragraph{Side channels of LLM services.} Encrypted streaming responses leak per-token lengths~\cite{weiss2024keylog} and packet-size and timing sequences that reveal conversation topics~\cite{mcdonald2025whisper}; deployed mitigations add random padding to streamed messages~\cite{cloudflare2024padding,openai2026obfuscation}, and later work recovers content despite them~\cite{zhang2025netecho,li2026merged}. Output token count, observed through timing, reveals an LLM classifier's class and a translator's target language; \citet{zhang2024timewilltell} propose, among other defenses, padding a call's output to the longest expected length. Efficiency mechanisms add further timing channels: speculative decoding~\cite{wei2024speculation,carlini2024timing} and shared prompt or key--value caches~\cite{gu2025caching,wu2025kvcache}. These works concern a single model call and a network or timing observer. We study workflows, where the leak is in which calls happen, and the usage-record observer, which sees exact counts even when transport is protected; our pad-only and random-padding baselines generalize the single-call defenses above and are shown to be insufficient for workflows.

\paragraph{Metadata leakage in agents and routers.} Traffic analysis fingerprints agents and their tool-use workflows~\cite{zhang2025agentprint} and multi-agent tasks on cloud LLM services~\cite{pouryousef2026benchmark}, whose authors conclude that effective defenses must shape the whole traffic profile. \citet{dangol2026graph} show that communication-graph metadata of agent protocols leaks task class and that message counts and sequence shape remain a residual channel under padding, closed by constant-rate cover traffic at the transport layer. Most closely, \citet{chen2026routing} measures, on 1.7 million real requests, how logged LLM router decisions leak sensitive topics even with content logging disabled, surveys the logging defaults of eleven gateways and platforms, and evaluates parity post-processing of the router's decisions. Our work is complementary in three ways: it addresses whole multi-call workflows rather than a single routing decision, it preserves the workflow's decisions rather than changing them, and it provides a compile-time transformation with a machine-checked guarantee. SICC~\cite{xu2026sicc} compiles multi-agent workflows so that the \emph{form} of messages exchanged between agents is independent of private state and proves a noninterference theorem for the message transcript; it does not address the number, model or token counts of the underlying LLM calls, which is our channel. The two compose.

\paragraph{Information-flow control for agents.} CaMeL~\cite{debenedetti2025camel} and FIDES~\cite{costa2025fides} use capability and taint tracking to stop prompt injection and explicit exfiltration; both explicitly leave control-flow side channels open, which is the gap we address for the usage observer. Language-based agent control~\cite{zhou2026lbac} lists cost and resource types as a possible policy kind, and Agenra~\cite{shivakumar2026agenra} gives a Lean-verified core for agent authorization; our validator follows the same methodology of a verified core with explicit assumptions, for a confidentiality hyperproperty. Static cost analysis of agent workflows~\cite{dickerson2026loom,hou2026ial,khan2026budgets} bounds or catalogs spending, which is a unary property; we need relational bounds whose equality across secret branches the shaping enforces. LLM programming systems such as LMQL~\cite{beurer2023lmql}, DSPy~\cite{khattab2024dspy} and SGLang~\cite{zheng2024sglang} structure prompts and decoding but make no confidentiality claims.

\paragraph{Constant-time and constant-resource programming.} The idea of transforming secret branches so that both perform the same observable operations goes back to cross-copying~\cite{agat2000transforming} and the program-counter security model~\cite{molnar2005pc}; FaCT~\cite{cauligi2019fact} is a DSL whose compiler rewrites secret control flow into constant-time code, and \citet{ngo2017constant} combine information flow with amortized resource analysis to verify and synthesize constant-resource programs. Raccoon~\cite{rane2015raccoon} executes decoy paths, and secure multi-execution~\cite{devriese2010sme} runs a program once per security level. Relational cost analysis~\cite{cicek2017relcost} and relational verification of side channels~\cite{antonopoulos2017blazer,chen2017themis} bound cost differences between runs, and verified compilers preserve constant-time and constant-resource properties~\cite{barthe2020ctcompcert,barthe2021constantresource}. \sys carries this line of work to a new resource model---LLM calls whose cost is set by a tokenizer and stochastic generation, observed per call by telemetry---with padding mechanisms specific to it (tokenizer-exact filler, forced decoding, model-matched dummy calls) and with translation validation~\cite{pnueli1998tv,necula2000tv} rather than a verified compiler.

\paragraph{Traffic-analysis defenses.} Website-fingerprinting defenses pad traffic to a fixed rate and size (BuFLO~\cite{dyer2012peekaboo}, Tamaraw~\cite{cai2014tamaraw}) or, more cheaply, to supersequences over anonymity sets~\cite{wang2014supersequence}; Pacer~\cite{mehta2022pacer} shapes cloud tenants' traffic so its shape is secret-independent by construction, and Shrinkwrap~\cite{bater2018shrinkwrap} pads intermediate result sizes of query plans with differential privacy. Our uniform strategy corresponds to fixed-size cells and our alignment to anonymity-set supersequences, applied to the structure of a workflow rather than to packets.
